% Genealogía operatoria añadida en REV.2.
% Residencia única del principio de contracción de fibras y memoria.

\chapter[Aritmética, cálculo discreto y memoria genealógica]{Genealogía
operatoria de la aritmética y el cálculo discreto: composición, contracción
de fibras y levantamientos con memoria}
\label{chap:genealogia-operaciones-rev2}
\index{operación!genealogía}
\index{memoria!levantamiento operatorio}

La aritmética publica resultados; una teoría genealógica debe conservar
además las condiciones que hicieron posible cada resultado. Esta distinción
no añade una operación nueva a la suma, el producto, la potencia, la raíz, el
logaritmo, la diferenciación o la integración. Precisa qué estructura
contrae cada una y qué datos debe transportar un levantamiento para distinguir
antecedentes que una misma salida identifica. APP exhibe el primer caso sobre
un soporte común; TRIT orienta localmente la transición; TPK conserva la
procedencia a lo largo de la composición.

\section{Operación publicada y levantamiento genealógico}

Sean \(X\) e \(Y\) conjuntos y sea \(f:X\to Y\) una aplicación que publica
una operación. Una \emph{memoria de operación} es un conjunto \(M_f\) junto
con una aplicación

\begin{equation}
 \widetilde f:X\longrightarrow Y\times M_f,
 \qquad
 \operatorname{pr}_Y\circ\widetilde f=f.
 \label{eq:levantamiento-genealogico-operacion-rev2}
\end{equation}

El primer componente publica la salida; el segundo registra los datos que el
lector decide conservar. El levantamiento es \emph{completo sobre las fibras}
si, para todo \(y\in Y\), la restricción

\[
 \operatorname{pr}_{M_f}\circ\widetilde f\big|_{f^{-1}(y)}
 :f^{-1}(y)\longrightarrow M_f
\]

es inyectiva. Esta condición expresa de manera exacta que dos antecedentes
con una misma salida visible permanecen distinguibles por su memoria. En las
realizaciones HMT no se exige que una única ficha conserve toda la historia:
la memoria completa puede distribuirse entre hoja, cociente, acarreo, fase,
orientación, frontera y registro cronológico enriquecido.

Si \(f:X\to Y\) y \(g:Y\to Z\) poseen levantamientos fijados, su composición
admite el levantamiento canónico asociado a esos levantamientos,

\begin{equation}
 \widetilde{g\circ f}:X\longrightarrow Z\times M_f\times M_g,
 \qquad
 \widetilde{g\circ f}(x)
 =\bigl(g(f(x)),m_f(x),m_g(f(x))\bigr),
 \label{eq:composicion-memorias-operacion-rev2}
\end{equation}

donde \(m_f\) y \(m_g\) son las componentes de memoria. La concatenación no
borra la procedencia de la primera etapa para sustituirla por la segunda: las
coordina. Éste es el principio operatorio que TPK realiza sobre rutas.

\begin{lemma}[Composición de levantamientos completos]
\label{lem:composicion-levantamientos-completos-rev2}
Si \(\widetilde f\) y \(\widetilde g\) son completos sobre las fibras de
\(f\) y \(g\), respectivamente, entonces
\(\widetilde{g\circ f}\) es completo sobre las fibras de \(g\circ f\).
\end{lemma}

\begin{proof}
Sean \(x_1,x_2\in X\) con
\(g(f(x_1))=g(f(x_2))\) y con iguales componentes de memoria en
\eqref{eq:composicion-memorias-operacion-rev2}. La igualdad de las memorias
\(m_g(f(x_1))=m_g(f(x_2))\), junto con la completitud de
\(\widetilde g\) sobre esa fibra de \(g\), implica
\(f(x_1)=f(x_2)\). La igualdad
\(m_f(x_1)=m_f(x_2)\) y la completitud de \(\widetilde f\) implican entonces
\(x_1=x_2\).
\end{proof}

\section{Suma y producto: alternativas y composición}

En un álgebra de caminos, la suma reúne alternativas con extremos
compatibles,

\[
 a_1\gamma_1+a_2\gamma_2,
\]

mientras el producto concatena etapas cuando el extremo de la primera
coincide con el origen de la segunda,

\[
 \gamma_2\gamma_1,
 \qquad t(\gamma_1)=s(\gamma_2).
\]

La suma conserva la coexistencia lineal de historias; el producto construye
una historia compuesta y, al publicarse sólo su valor, puede olvidar la
factorización que la produjo. Sobre la carta APP, las evaluaciones
\(\Sigma\) y \(\Pi\) actúan en el mismo estado \((i,j)\), pero organizan fibras
distintas. En el bloque central \(B_c\) de APP, la separación espectral que
se desarrolla en el capítulo siguiente mide exactamente esa incompatibilidad
local; no es una discrepancia entre dos tablas independientes ni se presenta
como una distancia global entre todos los lectores.

En la hoja aditiva, la traslación \(T_a(x)=x+a\) es reversible sobre el grupo
residual. En la hoja multiplicativa, \(M_a(x)=ax\) es invertible en
\(\mathbb Z/9\mathbb Z\) exactamente cuando \(\gcd(a,9)=1\). Las clases no
unitarias \(\overline0,\overline3,\overline6\), publicadas en la carta
positiva como \(9,3,6\), pliegan fibras y hacen visible la necesidad de
conservar cociente y procedencia. Suma y producto constituyen así los dos
lectores inaugurales de una misma geometría de caminos.

\begin{lemma}[Conjugación afín en el sector unitario]
\label{lem:conjugacion-afin-app-rev3}
Sea \(R=\mathbb Z/9\mathbb Z\). Para \(a\in R^\times\) y \(b\in R\), las
aplicaciones \(T_b(x)=x+b\) y \(M_a(x)=ax\) satisfacen
\begin{equation}
 M_aT_bM_a^{-1}=T_{ab}.
 \label{eq:conjugacion-afin-app-rev3}
\end{equation}
Si \(a\in(3)\), \(M_a\) deja de ser invertible y la conjugación no está
definida: su imagen queda contenida en el radical y varias ramas de la hoja
aditiva se identifican.
\end{lemma}

\begin{proof}
Para \(a\in R^\times\),
\(
 M_aT_bM_a^{-1}(x)=a(a^{-1}x+b)=x+ab=T_{ab}(x)
\).
Si \(a\in(3)\), el núcleo de \(M_a\) es no trivial; por tanto no existe
\(M_a^{-1}\) y la expresión de conjugación carece de tipo. El caso
\(a=3\) realiza explícitamente el plegado mediante
\(\ker M_3=\operatorname{im}M_3=(3)\).
\end{proof}

\section{Potencia, logaritmo y raíz}

Sea \(x\) un elemento de un monoide, o de un álgebra asociativa, y sea
\(n\geq1\). La potencia registra iteración de una composición:

\[
 x^n=\underbrace{x\cdots x}_{n\ \mathrm{factores}}.
\]

El exponente es la profundidad de esa iteración y la órbita de \(x\) determina
qué salidas se identifican. En el dominio real positivo, si \(b>0\),
\(b\neq1\) y \(x>0\), se cumple
\[
 \log_b(x^n)=n\log_b x.
\]
En el dominio complejo, la misma identidad requiere una rama compatible y
que la iteración no atraviese su corte. El logaritmo coordina la profundidad
\(n\) sólo después de fijar además el generador \(x\) y de exigir
\(\log_b x\neq0\). Base, generador, dominio, rama y corte forman parte del
lector.

La raíz invierte una potencia como problema de fibras,

\[
 y^n=x.
\]

La fibra puede ser vacía, unitaria o múltiple. Seleccionar una raíz requiere
datos de rama, orientación, fase o frontera. La construcción posterior de un
operador interno \(J_E\in\operatorname{End}(V)\), con \(J_E^2=-I_V\),
ejemplifica este principio: el
símbolo \(\sqrt{-1}\) publica una coordenada, mientras el operador, su dominio
y su acción conservan la estructura que la sostiene.

\section{Cálculo discreto sobre cilindros compatibles: operadores de diferenciación e integración}

Sea \(P_n\) un camino orientado conexo y sea \(A\) un grupo abeliano. En el
complejo discreto de medida, el coborde combinatorio

\[
 d:C^0(P_n;A)\longrightarrow C^1(P_n;A),
 \qquad (df)(e)=f(t(e))-f(s(e)),
\]

tiene por núcleo exactamente las \(0\)-cocadenas constantes. Reconstruir
\(f\) desde \(df\) exige una condición de borde, por ejemplo el valor en un
vértice de referencia.

El coborde depende de la incidencia y de la orientación, pero no de una
longitud. Si \(A\) es además un espacio vectorial real y
\(\ell:E(P_n)\to\mathbb R_{>0}\) es una longitud de arista, la derivada
métrica se define por
\[
 D_\ell f(e)=\frac{(df)(e)}{\ell(e)}.
\]
Para una ruta \(\gamma=(e_1,\ldots,e_r)\) compatible con las orientaciones,
\begin{equation}
 \sum_{k=1}^r D_\ell f(e_k)\,\ell(e_k)
 =f(t(\gamma))-f(s(\gamma)).
 \label{eq:integracion-derivada-metrica-rev2}
\end{equation}
La integración acumula así una cocadena a lo largo de una ruta y recupera la
diferencia de extremos cuando la cocadena es exacta. Coborde, derivada métrica
e integral dependen de datos distintos: incidencia y orientación; longitud y
escala; camino y sección de borde.

Esta dependencia explica por qué el paso de contar marcas a medir intervalos
es estructural. La derivación vive sobre relaciones intersticiales; la
integración compone esas relaciones; el ciclo añade el portador topológico de
la holonomía; el transporte TPK determina qué memoria sobrevive a la vuelta.

\section{Jerarquía de fibras y función del TPK}

La jerarquía puede condensarse sin identificar sus niveles:

\begin{center}
\small
\begin{tabularx}{\textwidth}{@{}
  >{\raggedright\arraybackslash}p{0.13\textwidth}
  >{\raggedright\arraybackslash}p{0.22\textwidth}
  >{\raggedright\arraybackslash}X
  >{\raggedright\arraybackslash}p{0.25\textwidth}@{}}
\toprule
\textbf{Operación} & \textbf{Acción primaria} &
\textbf{Fibra contraída por la salida} & \textbf{Memoria de reconstrucción}\\
\midrule
Suma & agregación o traslación & alternativas con una misma evaluación &
origen, coeficientes y soporte\\
Producto & composición o escala & factorizaciones y clases no unitarias &
orden de factores, cociente y hoja\\
Potencia & iteración & órbitas de composición &
exponente, orden y punto inicial\\
Raíz & inversión multivaluada & antecedentes de una potencia &
rama, fase, orientación y frontera\\
Logaritmo & coordenada de profundidad & periodos y determinaciones &
base, generador, dominio, rama y corte\\
Coborde y derivada métrica & diferencia local y normalización & modo constante &
incidencia, orientación, longitud, escala y referencia\\
Integral & acumulación de camino & rutas con igual valor publicado &
camino, medida, borde y holonomía\\
\bottomrule
\end{tabularx}
\end{center}

TPK conserva estas coordenadas dentro de estados transportables y no como una
lista añadida al resultado. Su ficha \((D,C,f,I,L,M)\) declara para cada
operador el dominio, el codominio, la regla, los invariantes, la información
publicada o perdida y la memoria conservada. La transición

\[
 \text{operación visible}
 \longrightarrow
 \text{levantamiento con memoria}
 \longrightarrow
 \text{sección compatible}
\]

proporciona el enlace entre la aritmética inaugural, la reconstrucción del
continuo y la definición posterior de número genealógico. La frontera
número--partícula conservará la misma arquitectura al añadir una realización
dimensional a la sección superviviente.

\paragraph{Dominio inaugural de la célula nonádica y exclusión causal de APP, TRIT, TPK y constantes analíticas}
La prioridad de suma y producto en APP, la conservación de fase, hoja,
acarreo y frontera por TPK, y la transición de sección a ruta constituyen las
premisas del levantamiento
\eqref{eq:levantamiento-genealogico-operacion-rev2}, la noción de completitud
sobre fibras y la tabla de memorias. No se identifica una operación con su realización física:
cada paso posterior conserva su dominio, su carta y su fuente.
